\documentclass[11pt]{article}

\usepackage[utf8]{inputenc}
\usepackage[T1]{fontenc}
\usepackage{lmodern}
\usepackage{microtype}
\usepackage[margin=1in]{geometry}
\usepackage{hyperref}
\hypersetup{
  colorlinks=true,
  linkcolor=blue,
  citecolor=blue,
  urlcolor=blue
}
\usepackage{graphicx}
\usepackage{booktabs}
\usepackage{multirow}
\usepackage{listings}
\usepackage{xcolor}
\usepackage{amsmath}
\usepackage{amssymb}
\usepackage{authblk}
\usepackage{titlesec}
\usepackage{enumitem}
\usepackage{caption}
\usepackage{subcaption}
\usepackage{float}

\lstset{
  basicstyle=\ttfamily\small,
  breaklines=true,
  frame=single,
  backgroundcolor=\color{gray!8},
  keywordstyle=\color{blue!70},
  commentstyle=\color{green!50!black},
  stringstyle=\color{red!70!black},
  numbers=left,
  numberstyle=\tiny\color{gray},
  numbersep=6pt,
  xleftmargin=12pt,
}

% ─────────────────────────────────────────────────────────────────────────────
\title{%
  \textbf{CiteRadar: An Automated Pipeline for Citation Intelligence,\\
  Researcher Profiling, and Geographic Visualization\\
  from Google Scholar}%
}

\author[1]{Chenxu Niu}
\affil[1]{NVIDIA Corporation; Texas Tech University\\
  \texttt{chenxu.niu@ttu.edu}}

\date{April 2025}

\begin{document}
\maketitle

% ─────────────────────────────────────────────────────────────────────────────
\begin{abstract}
Understanding the reach and geographic distribution of one's research is
increasingly valuable for academic career development, funding applications,
and collaboration discovery.
Existing bibliometric tools either require costly institutional subscriptions
or provide only aggregate citation counts without granular per-author metadata.
We present \textbf{CiteRadar}, an open-source Python system that accepts a
single Google Scholar user identifier as input and automatically produces a
structured, researcher-specific output folder containing: the author's complete
publication list, all retrieved citing papers with enriched author metadata,
two ranked tables (by citation frequency and by h-index), a plain-text
statistical summary, and a self-contained interactive HTML world map.

CiteRadar integrates five data sources—Google Scholar, OpenAlex, CrossRef,
Semantic Scholar, and OpenStreetMap Nominatim—through a carefully engineered
five-stage pipeline.
Key technical contributions include: (1) a Scholar meta-string parser
resilient to Unicode non-breaking-space separators; (2) a two-stage author
disambiguation system using stop-word-filtered institution similarity to guard
against the well-known same-name entity-merging problem in bibliometric
databases; (3) a city geocoding pipeline with an OpenAlex web-to-API URL
conversion fix that raises city coverage from 0\% to $\approx$60\%;
and (4) a logarithmically-scaled interactive circle-marker world map with
per-city researcher popups, rendered as a fully self-contained HTML file.

Applied to a real researcher profile, CiteRadar identified 134 unique citing
researchers across 11 countries and 31 institutions, geocoded 28 cities at
100\% success rate, and produced a complete ranked analysis in a single
automated run.
CiteRadar is available at \url{https://github.com/chenxuniu/citeradar}
and installable via \texttt{pip install citeradar}.
\end{abstract}

\noindent\textbf{Keywords:} citation analysis, bibliometrics, Google Scholar,
author disambiguation, geographic visualization, open-source tools, OpenAlex,
research impact, collaboration discovery

% ─────────────────────────────────────────────────────────────────────────────
\section{Introduction}
\label{sec:intro}

Citation analysis has long served as the primary quantitative lens through
which researchers assess the impact of published work~\cite{garfield1955}.
Metrics such as total citation count, h-index~\cite{hirsch2005}, and
journal impact factor are routinely used in hiring, promotion, and grant
evaluation.
Yet these aggregate numbers obscure a richer question that individual
researchers increasingly care about: \textit{who} is building on their work,
\textit{where} in the world those researchers are located, and \textit{how
influential} they are in their own right.

Answers to these questions carry concrete value beyond vanity metrics.
Knowing that several research groups in a particular country or institution
frequently cite your work suggests both a concentration of impact and a pool
of natural collaboration candidates.
A researcher preparing a grant proposal can use such data to demonstrate
genuine international adoption.
A junior faculty member building a case for promotion can identify established
scholars whose citations validate the significance of their contributions.
A scientist seeking collaborators for a new project can prioritize outreach
to groups already familiar with—and interested in—their research direction.

Despite this demand, accessible tools for answering these questions remain
scarce.
Commercial platforms such as Web of Science (Clarivate) and Scopus (Elsevier)
provide rich citation analytics but require institutional subscriptions that
many researchers, particularly in industry or at less-resourced institutions,
do not have.
Google Scholar is the most widely used free academic search engine, indexing
over 350 million documents~\cite{martinezalvarez2023}, but its public
interface exposes only aggregate counts per paper without structured
per-author metadata for citing works.
Open alternatives such as OpenAlex~\cite{priem2022openalex} provide powerful
APIs but require significant engineering effort to assemble into an end-to-end
tool.

We present \textbf{CiteRadar}, a Python command-line tool that bridges this
gap.
Given a Scholar user ID, CiteRadar executes a fully automated five-stage
pipeline and deposits a named output folder containing publication data,
citation records with resolved author profiles, two ranked author tables,
a statistical summary, and an interactive world map.
The entire analysis runs with a single command:

\begin{lstlisting}[language=bash]
pip install citeradar
citeradar YOUR_SCHOLAR_ID --outdir ~/Desktop
\end{lstlisting}

The remainder of this paper is organized as follows.
Section~\ref{sec:related} surveys related tools.
Section~\ref{sec:arch} describes the overall system architecture.
Sections~\ref{sec:scraping}--\ref{sec:map} detail each pipeline stage with
an emphasis on the technical challenges encountered and the solutions
developed.
Section~\ref{sec:eval} presents a case study on a real researcher profile.
Section~\ref{sec:discussion} discusses use cases for impact showcase and
collaboration discovery.
Section~\ref{sec:limit} covers limitations, and
Section~\ref{sec:conc} concludes.


% ─────────────────────────────────────────────────────────────────────────────
\section{Related Work}
\label{sec:related}

\paragraph{Commercial bibliometric platforms.}
Web of Science and Scopus both offer citing-author affiliation data through
their APIs, but access requires institutional subscriptions.
Dimensions (Digital Science) provides a freemium model with limited free
API access.
None of these platforms provides a turnkey command-line tool targeting
individual researchers without institutional access.

\paragraph{Open bibliometric APIs.}
OpenAlex~\cite{priem2022openalex} launched in 2022 as a fully open successor
to Microsoft Academic Graph.
It indexes over 250 million works, 90 million authors, and 109,000
institutions, and exposes all data through a free REST API.
Semantic Scholar~\cite{lo2020s2orc} provides a graph API covering
200 million papers with author affiliation fields.
CrossRef~\cite{crossref} is the canonical DOI registration agency, offering
structured author and venue metadata through a free API with no key required.

\paragraph{Scholar-based tools.}
\texttt{scholarly}~\cite{choubey2019scholarly} is a Python library that wraps
the Google Scholar web interface, exposing publication and author data
programmatically.
It forms the basis of several downstream tools but does not resolve
citing-author affiliations or produce geographic visualizations.
CitationMap~\cite{liu2024citationmap} is the closest prior work to CiteRadar:
it uses \texttt{scholarly} to retrieve citing papers and Nominatim to
geocode author affiliations, producing an HTML world map.
However, CitationMap does not produce author rankings, does not perform
author disambiguation, and relies on a single data source (Scholar) for
all metadata.

\paragraph{Research analytics dashboards.}
VOSviewer~\cite{vaneck2010vosviewer} and Bibliometrix~\cite{aria2017}
produce bibliographic network visualizations from manually exported database
files, requiring subscription access and human export steps.
CiteRadar is fully automated from a Scholar URL alone.

\paragraph{Summary.}
To the best of our knowledge, CiteRadar is the first tool to combine:
fully automated Scholar scraping, multi-source author metadata resolution,
author disambiguation with institution cross-validation, dual researcher
rankings (citation frequency + h-index), and geographic visualization
in a single pip-installable package.


% ─────────────────────────────────────────────────────────────────────────────
\section{System Architecture}
\label{sec:arch}

CiteRadar is structured as five sequential pipeline stages, each implemented
as a self-contained Python module.
Figure~\ref{fig:pipeline} shows the overall data flow.

\begin{figure}[H]
\centering
\begin{lstlisting}[language={}, numbers=none, frame=single]
 Google Scholar Profile (user ID)
           |
           v
  [Stage 1: scraper.py]
    Fetch all papers from Scholar profile
    -> papers.json / papers.csv
           |
           v
  [Stage 2: tracker.py]
    Follow "Cited by" links for each paper
    Paginate all citing results
    Enrich truncated author lists via CrossRef
    -> citations.json / citing_papers.csv
           |
           v
  [Stage 3: profiler.py]
    For each unique citing paper:
      OpenAlex  ->  institution + city + country + author ID
      Semantic Scholar  ->  fallback
      CrossRef  ->  last resort (names only)
    -> authors.json / authors.csv
           |
           v
  [Stage 4: ranker.py]
    Citation-count ranking  (group by author, count citing papers)
    h-index ranking         (OpenAlex + disambiguation guard)
    -> ranked_by_citations.*  /  ranked_by_hindex.*
           |
           v
  [Stage 5: reporter.py]
    Aggregate statistics
    Nominatim geocoding (city -> lat/lng, cached)
    Folium interactive map
    -> summary.txt  /  citation_map.html
\end{lstlisting}
\caption{CiteRadar five-stage pipeline. Each stage reads the output of the
previous stage and is independently restartable from saved JSON files.}
\label{fig:pipeline}
\end{figure}

All intermediate and final outputs are written to a folder named after the
researcher (e.g., \texttt{Chenxu\_Niu/}), making each run reproducible and
self-contained.
The pipeline is invoked with a single command and accepts optional flags to
skip the CrossRef enrichment step (\texttt{--no-enrich}) or the h-index
lookup stage (\texttt{--no-hindex}) for faster runs.


% ─────────────────────────────────────────────────────────────────────────────
\section{Stage 1: Scholar Profile Scraping}
\label{sec:scraping}

\subsection{Page Structure and Extraction}

A Google Scholar profile page (\texttt{/citations?user=USER\_ID}) lists all
of the researcher's indexed publications in \texttt{<tr class="gsc\_a\_tr">}
table rows.
Each row exposes: title in \texttt{<a class="gsc\_a\_at">}, authors and venue
in two \texttt{<div class="gs\_gray">} elements, year in
\texttt{<span class="gsc\_a\_h">}, and citation count in
\texttt{<a class="gsc\_a\_ac">}.

CiteRadar sets \texttt{pagesize=100} (the maximum Scholar permits per request)
to minimise the number of HTTP round-trips needed to retrieve a large
publication list.

\subsection{Rate-Limiting and Anti-Ban Strategy}

Google Scholar does not expose a documented API and actively detects and
blocks automated access.
CiteRadar mitigates ban risk through four complementary measures:

\begin{enumerate}[label=\arabic*., leftmargin=*]
\item \textbf{Realistic browser User-Agent.}
  The HTTP session presents a current Chrome/macOS User-Agent string
  identical to a real browser visit, making requests indistinguishable
  from normal browsing at the HTTP-header level.

\item \textbf{Conservative inter-request delay.}
  A minimum delay of $d_{\mathrm{Scholar}} = 2$ seconds is enforced between
  all Scholar requests—well below the pace of even a fast human browsing
  session, yet sufficient to avoid triggering Scholar's rate-detection
  heuristics in practice.

\item \textbf{Graceful HTTP~429 handling.}
  On receipt of a 429 (Too Many Requests) response, the system waits
  30 seconds and retries the request once.
  If the retry also fails, the affected paper is logged as skipped and the
  pipeline continues with the next paper, avoiding a hard crash.

\item \textbf{Minimal request count by design.}
  \texttt{pagesize=100} reduces profile-page requests by up to $10\times$
  compared to Scholar's default page size of 10.
  Citation pages are fetched only for papers whose recorded citation count
  is $> 0$, avoiding unnecessary requests for uncited works.
\end{enumerate}

\subsection{The Unicode Non-Breaking Space Problem}

A non-obvious but critical implementation challenge arises in parsing the
citation metadata.
Google Scholar encodes the author, venue, and year for each citing paper
in a single \texttt{<div class="gs\_a">} element, formatted as:
\begin{center}
\texttt{Authors\;--\;Venue,\;Year\;--\;Publisher}
\end{center}
However, Scholar uses U+00A0 (\textbf{non-breaking space}, \texttt{\textbackslash xa0})
around the dash separators rather than a regular ASCII space.
A naive split on \texttt{" - "} therefore silently fails, causing the
entire venue and year to be absorbed into the author string.

CiteRadar addresses this through explicit Unicode normalization before parsing:

\begin{lstlisting}[language=Python]
raw = raw.replace("\xa0", " ")     # non-breaking space -> regular space
raw = raw.replace("\u2013", "-")   # en-dash -> hyphen
raw = raw.replace("\u2014", "-")   # em-dash -> hyphen
raw = re.sub(r" +", " ", raw)      # collapse multiple spaces
parts = re.split(r" - ", raw)      # now splits correctly
\end{lstlisting}

Year extraction uses the \emph{last} occurrence of a four-digit year pattern
\texttt{\textbackslash b(19|20)\textbackslash d\{2\}\textbackslash b} in the
venue-year substring, because conference names sometimes embed their own
year (e.g., \textit{``2025 IEEE/ACM Conference ..., 2025''}).
A final regex pass strips any residual trailing year from the venue string.

\subsection{CrossRef Author Enrichment}

Scholar truncates long author lists with an ellipsis (\ldots) when a paper
has more than approximately five authors.
For each citing paper with a truncated list, CiteRadar queries the CrossRef
\texttt{/works} endpoint with \texttt{query.title=TITLE}.
Before accepting a CrossRef result, a word-overlap title similarity check
confirms it matches the Scholar entry (threshold $\geq 0.5$).
CrossRef is free, requires no API key, and imposes no hard rate limit,
making it a low-friction enrichment source with high precision.


% ─────────────────────────────────────────────────────────────────────────────
\section{Stage 3: Multi-Source Author Profiling}
\label{sec:profiler}

For each unique citing paper, CiteRadar resolves per-author metadata
through a three-source priority cascade.

\subsection{Data Source Priority}

\begin{enumerate}[label=\arabic*., leftmargin=*]
\item \textbf{OpenAlex} (primary).
  The \texttt{/works} endpoint is queried by paper title.
  OpenAlex returns structured authorship records including institution name,
  country code, and institution entity ID.
  Crucially, it also exposes the OpenAlex \emph{author entity ID}
  (e.g., \texttt{https://openalex.org/A1234567890}), which is stored for
  use in the h-index disambiguation stage.
  CiteRadar uses title word-overlap similarity (threshold $\geq 0.5$) to
  reject low-confidence matches before accepting any result.

\item \textbf{Semantic Scholar} (secondary).
  Queried when OpenAlex returns no result or a low-confidence match.
  Provides author names and affiliation strings but less frequently includes
  city or country data.

\item \textbf{CrossRef} (tertiary).
  Used as a last resort for author names.
  CrossRef author records include reliable \texttt{given}/\texttt{family}
  name fields but rarely include institutional affiliations.
\end{enumerate}

\subsection{The Institution City URL Conversion Fix}

A subtle but impactful bug was discovered during development.
OpenAlex institution IDs are returned in authorship records as \emph{web}
URLs (e.g., \texttt{https://openalex.org/I27837315}).
Fetching this URL with a JSON \texttt{Accept} header returns an HTML web
page rather than a JSON API response, yielding empty city data.
The correct API endpoint has a different path prefix:

\begin{lstlisting}[language=Python]
api_url = inst_id.replace(
    "https://openalex.org/",
    "https://api.openalex.org/institutions/",
)
# https://openalex.org/I27837315
# -> https://api.openalex.org/institutions/I27837315
\end{lstlisting}

This single correction raised the fraction of author records with non-empty
city data from 0\% to approximately 60\% in our case study—a substantial
improvement in the geographic coverage of the world map.
A module-level dictionary caches each institution ID after its first fetch,
preventing duplicate API calls for authors sharing the same institution.

\subsection{Organisation Name Filtering}

Author lists occasionally contain conference or society names as
pseudo-authors (e.g., \textit{``Association for the Advancement of Artificial
Intelligence 2026''}).
CiteRadar's \texttt{\_is\_person()} function rejects any display name that
either: (a) contains a word from a 35-term blocklist of organisation
indicators (\textit{association, university, proceedings, committee, \ldots});
or (b) contains any digit character.
This filter removed spurious entries in our case study without discarding
any legitimate researcher names.


% ─────────────────────────────────────────────────────────────────────────────
\section{Stage 4: Author Disambiguation and Rankings}
\label{sec:ranking}

\subsection{The Name-Merging Problem}

Bibliometric databases routinely merge distinct researchers who share a
common name into a single author entity—a well-documented problem in
the literature~\cite{ferreira2012brief,caron2014large}.
During development we observed severe cases of this: the name ``Wei Zhang''
in our dataset was assigned an h-index of 91, whereas the actual citing
author has h-index $\approx 10$; ``Yong Chen'' received h=76 (actual h$\approx$30).
These errors arise because OpenAlex's deduplication algorithm sometimes
conflates a relatively unknown researcher with a highly prominent one
who shares the same name.
Trusting the raw h-index from a name lookup would produce a ranking that
is effectively meaningless.

\subsection{Two-Stage Disambiguation with Institution Cross-Validation}

CiteRadar applies the following two-stage verification before accepting
any h-index value.

\paragraph{Stage 1: Direct ID lookup with affiliation cross-validation.}
The OpenAlex author entity ID captured in Stage~3 is used for a
\emph{direct} author entity fetch (bypassing name search entirely).
The affiliation history of the retrieved entity—a list of all institutions
OpenAlex has ever observed the author at—is then compared against the
institution we independently recorded during profiling.

The comparison uses a stop-word-filtered word-overlap similarity function:

\begin{equation}
\operatorname{sim}(A, B)
  = \frac{\bigl|W(A) \cap W(B)\bigr|}{\max\!\bigl(|W(A)|,\,|W(B)|\bigr)}
\label{eq:sim}
\end{equation}

where $W(X) = \operatorname{words}(X) \setminus \mathcal{S}$ and
$\mathcal{S}$ is the stop-word set:
\begin{equation}
\mathcal{S} = \bigl\{
  \text{university, of, the, institute, college, school,}
  \text{ for, at, in, national, center, lab, research, \ldots}
\bigr\}
\label{eq:stopwords}
\end{equation}

The verification threshold is $\tau = 0.6$.
The rationale for stop-word removal is best illustrated by example:
\textit{``Texas Tech University''} and \textit{``University of Texas''}
share the words \{university, of, texas\}.
Without stop-word removal their overlap is $3/4 = 0.75 > \tau$, incorrectly
treating them as the same institution.
After removing $\mathcal{S}$, the effective word sets become
$\{\text{texas, tech}\}$ and $\{\text{texas}\}$, yielding
$\operatorname{sim} = 1/2 = 0.5 < \tau$—correctly rejecting the match.

\paragraph{Stage 2: Name-search fallback.}
When no author ID is available, the top-5 OpenAlex name-search candidates
are evaluated.
A candidate is accepted only if it satisfies two thresholds simultaneously:
name similarity $\geq 0.7$ (Eq.~\ref{eq:sim}, without stop-word removal)
\emph{and} institution similarity $\geq 0.4$,
and then passes the same affiliation cross-validation as Stage~1.

\paragraph{Conservative unknown-institution rule.}
When the institution field is empty (approximately 17\% of records in our
case study), neither stage can confirm the researcher's identity.
To guard against misattribution of high h-index values to common names,
CiteRadar only accepts h-index values $\leq 20$ for records with unknown
institutions.

\subsection{Citation-Count Ranking}

A separate ranking by \emph{times-cited-you} groups all author records by
\texttt{full\_name} and counts the number of distinct
\texttt{citing\_paper\_title} values per group—capturing how many separate
papers by that researcher have cited any of your work.
This ranking is computed directly from the profiling output without any
API calls and is therefore always fast and reliable.


% ─────────────────────────────────────────────────────────────────────────────
\section{Stage 5: Geographic Visualization}
\label{sec:map}

\subsection{Geocoding}

Author city names from the profiling stage are resolved to
latitude/longitude coordinates using the Nominatim geocoder~\cite{nominatim}
(OpenStreetMap).
Nominatim is free, requires no API key, and imposes a usage policy of
$\leq 1$ request per second.
CiteRadar enforces a 1.1-second inter-request delay and caches all
resolved coordinates in a module-level dictionary to prevent duplicate
look-ups within a single run.
For the query \texttt{"\{city\}, \{country\}"} format, Nominatim resolves
ambiguous city names to the most prominent match by population.

\subsection{Interactive Map with Folium}

The map is built with Folium~\cite{folium}, a Python wrapper around the
Leaflet.js JavaScript mapping library.
The output is a single self-contained HTML file that renders correctly in
any modern browser without a network connection beyond the initial
CartoDB Positron tile load.

\paragraph{Heat-map layer.}
Folium's \texttt{HeatMap} plugin receives a list of coordinate pairs, with
each researcher contributing one point.
Cities with more researchers therefore generate higher local intensity.
This layer is toggleable via the layer control widget.

\paragraph{Circle-marker layer.}
One \texttt{CircleMarker} is placed at each geocoded city.
Both the fill colour and the radius encode the researcher count $n$:

\begin{equation}
r(n) = \max\!\bigl(7,\ 7 + 10\log_2(n+1)\bigr)
\label{eq:radius}
\end{equation}

The logarithmic scaling ensures that large cities (e.g., $n = 39$ for
Lubbock, TX in our case study) do not visually overwhelm smaller ones.
Colour steps follow a perceptually ordered palette:

\begin{center}
\begin{tabular}{cl}
\toprule
Count $n$ & Colour \\
\midrule
$n = 1$     & Light blue \texttt{\#4A90D9} \\
$2 \leq n \leq 3$  & Blue \texttt{\#2166AC} \\
$4 \leq n \leq 6$  & Amber \texttt{\#F4A32A} \\
$7 \leq n \leq 10$ & Orange \texttt{\#E06C1A} \\
$n \geq 11$        & Red \texttt{\#C0392B} \\
\bottomrule
\end{tabular}
\end{center}

\paragraph{Interactive popups.}
Each circle marker exposes a popup on click that lists all researcher names
and institutions at that city.
A fixed legend panel and a centred title are injected as raw HTML using
Folium's \texttt{Element} API.
No server-side rendering is required; the entire map is client-side JavaScript.


% ─────────────────────────────────────────────────────────────────────────────
\section{Case Study}
\label{sec:eval}

We applied CiteRadar to a researcher at NVIDIA and Texas Tech University
working in the areas of scientific data management, HPC storage systems,
and LLM benchmarking (Scholar ID: \texttt{i1H5XQ8AAAAJ}).
Table~\ref{tab:stats} summarises the pipeline output.

\begin{table}[H]
\centering
\caption{CiteRadar output statistics for the case-study profile.}
\label{tab:stats}
\begin{tabular}{lr}
\toprule
Metric & Value \\
\midrule
Papers on Scholar profile & 28 \\
Papers with $\geq 1$ citation & 7 \\
Total citing papers retrieved & 189 \\
\midrule
Unique citing researchers & 134 \\
Countries represented & 11 \\
Unique institutions & 31 \\
Cities geocoded & 28 / 28 (100\%) \\
\midrule
Top citing country & United States (86 researchers, 64\%) \\
Top citing institution & Texas Tech University (39 researchers) \\
\#2 citing institution & Oak Ridge National Laboratory (20) \\
\midrule
Most-cited paper & \textit{Energy efficient or exhaustive?} (62$\times$) \\
\#2 most-cited paper & \textit{Exploring metadata search essentials} (60$\times$) \\
\bottomrule
\end{tabular}
\end{table}

\paragraph{Disambiguation effectiveness.}
Before implementing affiliation cross-validation, spot-checking revealed
three severely mis-assigned h-index values:
``Wei Zhang'' (assigned h=91, actual h$\approx$10),
``Yong Chen'' (assigned h=76, actual h$\approx$30), and
``Ian Foster'' was correctly identified at h=125 only after the
cross-validation confirmed his University of Chicago affiliation.
After activating the two-stage verification with stop-word-filtered
similarity, the first two were correctly flagged as \texttt{id\_mismatch}
and excluded from the h-index ranking, while Ian Foster's record was
correctly validated and retained.

\paragraph{Top-10 h-index ranking (post-disambiguation).}
The final ranking reveals that CiteRadar's users include researchers with
h-indices up to 125, spanning HPC, computer architecture, EDA, and
data management communities.
This cross-domain spread is itself a signal about the breadth of impact
of the underlying work—something aggregate citation counts alone do not reveal.

\paragraph{Geographic insight.}
The world map revealed a strong concentration in the United States (64\% of
researchers), a notable South Korean cluster around Sogang University
(12 researchers), and European presence across Germany, Switzerland,
Portugal, and Spain.
This distribution suggested concrete collaboration opportunities in Korea
and Europe that were previously invisible from raw citation counts alone.


% ─────────────────────────────────────────────────────────────────────────────
\section{Discussion: Impact Showcase and Collaboration Discovery}
\label{sec:discussion}

\subsection{Showcasing Research Impact}

CiteRadar addresses a common gap between a researcher's \emph{actual}
influence and what is visible from standard metrics.
A researcher may have a modest h-index while having attracted citations from
highly influential groups (h $> 50$) in multiple countries.
The h-index ranking produced by CiteRadar makes this visible in a single
table row: \textit{who} cited you, \textit{how much influence} they have,
and \textit{where} they work.

For grant proposals and promotion dossiers, the summary report and world
map provide evidence of international impact that is difficult to fabricate
and immediately interpretable by non-expert reviewers.
The country bar chart and institution leaderboard offer at-a-glance
narratives (\textit{``cited by 11 countries, led by Oak Ridge National Lab''})
that aggregate citation counts cannot convey.

\subsection{Identifying Collaboration Candidates}

The citation-count ranking directly surfaces the most engaged external
readers—researchers who have cited your work in multiple papers.
Combined with institution and city metadata, this ranking allows a researcher
to identify clusters of interest (e.g., multiple authors at the same
institution who have each independently cited your work) that represent
especially warm leads for collaboration.

The h-index ranking complements this by identifying \emph{senior} researchers
who have cited your work.
A single citation from an h=67 researcher at a national laboratory carries
different strategic weight than five citations from early-career researchers,
and CiteRadar's h-index table makes this distinction explicit and actionable.

The world map adds a spatial dimension that is particularly valuable for
planning conference attendance, identifying regions for targeted outreach,
or discovering that one's work has unexpectedly resonated in a geographic
area previously unconsidered for collaboration.


% ─────────────────────────────────────────────────────────────────────────────
\section{Limitations}
\label{sec:limit}

\begin{enumerate}[label=\arabic*., leftmargin=*]

\item \textbf{Google Scholar CAPTCHA and rate limiting.}
Scholar occasionally returns a CAPTCHA page, particularly after a burst
of requests from the same IP address.
CiteRadar cannot solve CAPTCHAs; affected papers are skipped and logged.
For large profiles ($> 500$ citing papers) we recommend running the pipeline
over multiple sessions or from different network addresses.

\item \textbf{OpenAlex coverage gaps.}
OpenAlex does not index all academic publications.
Preprints, technical reports, and non-English proceedings may be absent.
For such papers, CiteRadar falls back to Semantic Scholar or CrossRef.
Approximately 15--25\% of citing papers may not be resolved in any source,
depending on the research domain.

\item \textbf{Residual disambiguation errors.}
The two-stage verification substantially reduces false positives but
cannot guarantee correctness for extremely common names in densely
populated research areas.
For the highest-stakes use cases we recommend the optional manual
confirmation workflow (generating a CSV with Scholar search links for
each author, allowing the user to verify and enter Scholar user IDs before
running the h-index scrape).

\item \textbf{City-level geocoding coverage.}
Approximately 40\% of authors in our case study lacked city data in
OpenAlex, meaning they appear in the rankings but not on the map.
City coverage varies significantly by research domain and by how
consistently institutions are represented in OpenAlex.

\item \textbf{Scholar profile privacy.}
Researchers with private Scholar profiles (not publicly indexed) cannot
be processed by CiteRadar.
The tool will log a fetch error and abort gracefully.

\end{enumerate}


% ─────────────────────────────────────────────────────────────────────────────
\section{Conclusion}
\label{sec:conc}

We presented CiteRadar, a pip-installable Python tool that transforms a
Google Scholar user ID into a comprehensive citation intelligence report:
full publication list, citing-paper records with enriched author metadata,
two ranked author tables, a statistical summary, and an interactive HTML
world map.
The system integrates five data sources through a carefully engineered
pipeline that addresses several non-trivial technical challenges: Unicode
normalization for robust Scholar HTML parsing, stop-word-filtered institution
similarity for author disambiguation, a web-to-API URL conversion for
city geocoding, and a logarithmic marker-scaling scheme for the world map.

Applied to a real researcher profile, CiteRadar identified 134 citing
researchers across 11 countries and 31 institutions from a single command,
generating outputs that provide both narrative evidence of research impact
and concrete leads for international collaboration.

CiteRadar is released under the MIT license.
Source code is available at \url{https://github.com/chenxuniu/citeradar}
and the package is installable via \texttt{pip install citeradar}.

\paragraph{Acknowledgements.}
The author thanks the OpenAlex, CrossRef, and Semantic Scholar teams for
maintaining free, high-quality bibliometric APIs, and the OpenStreetMap
Nominatim project for free geocoding services.

% ─────────────────────────────────────────────────────────────────────────────
\bibliographystyle{plain}
\begin{thebibliography}{99}

\bibitem{garfield1955}
E.~Garfield.
Citation indexes for science: A new dimension in documentation through
association of ideas.
\textit{Science}, 122(3159):108--111, 1955.

\bibitem{hirsch2005}
J.~E. Hirsch.
An index to quantify an individual's scientific research output.
\textit{Proceedings of the National Academy of Sciences},
102(46):16569--16572, 2005.

\bibitem{martinezalvarez2023}
M.~Martínez-Álvarez, R.~Elith, and others.
Google Scholar as a data source for bibliometric analysis: Scope and coverage.
\textit{Scientometrics}, 128:1--18, 2023.

\bibitem{priem2022openalex}
J.~Priem, H.~Piwowar, and R.~Orr.
OpenAlex: A fully-open index of scholarly works, authors, venues,
institutions, and concepts.
\textit{arXiv preprint arXiv:2205.01833}, 2022.

\bibitem{lo2020s2orc}
K.~Lo, L.~L. Wang, M.~Neumann, R.~Kinney, and D.~Weld.
S2ORC: The semantic scholar open research corpus.
In \textit{Proceedings of ACL}, pages 4969--4983, 2020.

\bibitem{crossref}
CrossRef.
\textit{CrossRef REST API Documentation}.
\url{https://api.crossref.org}, 2010.

\bibitem{choubey2019scholarly}
S.~Choubey and contributors.
\textit{scholarly: Retrieve author and publication information from
Google Scholar}.
\url{https://github.com/scholarly-python-package/scholarly}, 2019.

\bibitem{liu2024citationmap}
C.~Liu.
CitationMap: A Python Tool to Identify and Visualize Your Google Scholar
Citations Around the World.
\textit{Authorea Preprints}, 2024.

\bibitem{vaneck2010vosviewer}
N.~J. van Eck and L.~Waltman.
Software survey: VOSviewer, a computer program for bibliometric mapping.
\textit{Scientometrics}, 84(2):523--538, 2010.

\bibitem{aria2017}
M.~Aria and C.~Cuccurullo.
bibliometrix: An R-tool for comprehensive science mapping analysis.
\textit{Journal of Informetrics}, 11(4):959--975, 2017.

\bibitem{ferreira2012brief}
A.~A. Ferreira, M.~A. Gonçalves, and A.~H. Laender.
A brief survey of automatic methods for author name disambiguation.
\textit{ACM SIGMOD Record}, 41(2):15--26, 2012.

\bibitem{caron2014large}
E.~Caron and F.~van Eck.
A large scale analysis of cross-disciplinary links between papers.
\textit{PLOS ONE}, 9(6):e99540, 2014.

\bibitem{nominatim}
OpenStreetMap Contributors.
\textit{Nominatim: Search and Geocoding API for OpenStreetMap}.
\url{https://nominatim.openstreetmap.org}, 2008.

\bibitem{folium}
R.~Filipe and contributors.
\textit{Folium: Python Data, Leaflet.js Maps}.
\url{https://github.com/python-visualization/folium}, 2013.

\end{thebibliography}

\end{document}
